% Pista geo-24j-q6 · Geometria
% Procedència: PAU Catalunya 2024 · Convocatòria de juny · Sèrie 1 · Qüestió 6
\documentclass[11pt,a4paper]{article}
\usepackage{pista}
\pistainfo{Catalunya 2024}{Convocatòria de juny}{Sèrie 1}

\begin{document}

\section*{\textcolor{blauPista}{Qüestió 6}}

\noindent Considereu els punts $A = (1, 2, 3)$ i $B = (-3, -2, 3)$.

\subsection*{\textit{a)} \normalfont Calculeu l'equació del pla $\pi$ que és perpendicular a la recta $AB$ i que passa pel punt mitjà entre $A$ i $B$. Justifiqueu que aquest pla està format, precisament, pels punts $P = (x, y, z)$ que estan a igual distància de $A$ que de $B$, és a dir, $d(P, A) = d(P, B)$. \hspace*{\fill}\textnormal{[1~punt]}}

\pista{Aquest apartat et demana dues coses.

\medskip
\textit{El pla $\pi$.} Com que el pla és perpendicular a la recta $AB$, el seu vector normal és precisament el vector $\vec{AB}$. En primer lloc, calcula $\vec{AB}$ i el punt mitjà entre $A$ i $B$. A continuació, escriu el pla en la forma \mbox{$ax + by + cz + d = 0$}, amb $(a, b, c) = \vec{AB}$, i troba $d$ fent que el pla passi pel punt mitjà.

\medskip
\textit{Justificació.} Escriu $d(P, A) = d(P, B)$ amb $P = (x, y, z)$. Per no haver de treballar amb arrels quadrades, eleva els dos costats al quadrat:
\[
d(P, A)^{2} = d(P, B)^{2}
\]
Desenvolupa els quadrats i simplifica. No és tan difícil com sembla, perquè els termes $x^{2}$, $y^{2}$ i $z^{2}$ s'anul·len. Comprovaràs que obtens la mateixa equació del pla $\pi$.}

\subsection*{\textit{b)} \normalfont Calculeu les distàncies de $A$ i de $B$ al pla $\pi$ i comproveu que són iguals. És casualitat? Raoneu la resposta. \hspace*{\fill}\textnormal{[0,75~punts]}}

\pista{Aquest apartat et demana dues coses.

\medskip
\textit{Les distàncies.} Aplica a $A$ i a $B$ la fórmula de la distància d'un punt $X = (x_{0}, y_{0}, z_{0})$ a un pla, amb el pla de l'apartat \textit{a)} en la forma \mbox{$ax + by + cz + d = 0$}:
\[
d(X, \pi) = \dfrac{|a x_{0} + b y_{0} + c z_{0} + d|}{\sqrt{a^{2} + b^{2} + c^{2}}}
\]

\medskip
\textit{És casualitat?} Recorda com has construït el pla $\pi$: és perpendicular a la recta $AB$ i passa pel punt mitjà de $A$ i $B$. Fes-ne un dibuix: el pla talla el segment $AB$ just pel mig, i formant un angle recte. Què pots dir, llavors, de les distàncies de $A$ i de $B$ al pla?}

\subsection*{\textit{c)} \normalfont Sigui $C = (-7, 6, 3)$. El triangle $ABC$ és isòsceles? Calculeu la seva àrea. \hspace*{\fill}\textnormal{[0,75~punts]}}

\pista{En primer lloc, comprova si el punt $C$ pertany al pla $\pi$ de l'apartat \textit{a)}. Si ho és, què en dedueixes sobre les distàncies de $C$ a $A$ i de $C$ a $B$? Recorda què has justificat a l'apartat \textit{a)}. En canvi, si no ho és, calcula les longituds dels tres costats i compara-les.

\medskip
A continuació, calcula l'àrea amb el costat $AB$ com a base. Si el triangle és isòsceles, l'altura sobre $AB$ va de $C$ al punt mitjà $M$ de $A$ i $B$. Per tant:
\[
\text{Àrea} = \frac{\text{base} \cdot \text{altura}}{2} = \frac{|\vec{AB}| \cdot d(C, M)}{2}
\]}

\end{document}
